%!TEX root = ../../note.tex
\subsection{Numbers and the Language of Logic}

\subsubsection{The number systems}
We start from the familiar systems
\begin{align*}
    \N &= \set{1,2,3,\ldots}, \\
    \Z &= \set{\ldots,-1,0,1,\ldots}, \\
    \Q &= \set{\frac{m}{n} : m,n\in\Z,\ n\ne 0}.
\end{align*}
Each one repairs a defect of the previous one. $\N$ is closed under
addition; $\Z$ is also closed under subtraction and multiplication, but not
under division ($1/2\notin\Z$); $\Q$ admits all quotients with nonzero
denominator. Even $\Q$, however, is too small: it has no number whose square
is $2$.

\begin{theorem}[Irrationality of $\sqrt2$]\label{thm:sqrt2-irrational}
    $\sqrt{2}\notin\Q$.
\end{theorem}

The idea is to write $\sqrt2$ in lowest terms; the equation $m^2=2n^2$ then
forces both $m$ and $n$ to be even.

\begin{proof}
    Suppose $\sqrt2=m/n$ with $m,n\in\Z$, $n>0$, $\gcd(m,n)=1$. Squaring
    gives $m^2=2n^2$, so $m^2$ is even. Since an odd integer has an odd
    square, $m$ is even, say $m=2\ell$. Then $4\ell^2=2n^2$, i.e.\
    $n^2=2\ell^2$, so $n$ is even by the same argument. Thus $2$ divides
    both $m$ and $n$, contradicting $\gcd(m,n)=1$.
\end{proof}

The elements of $\R\setminus\Q=\set{x\in\R : x\notin\Q}$ are called
\term{irrational numbers}.

\subsubsection{Decimal representation}
Informally, every nonnegative real number has a decimal expansion
\begin{equation*}
    x=a+\sum_{i=1}^{\infty}\frac{x_i}{10^i},
    \qquad a\in\N\cup\set{0},
    \qquad x_i\in\set{0,1,\ldots,9},
\end{equation*}
and negative numbers are negatives of such expansions. The expansion need
not be unique.

\begin{example}
    $1.999\cdots=2$. Indeed, if $x=1.999\cdots$, then $10x=19.999\cdots$, so
    $9x=10x-x=18$ and $x=2$.
\end{example}

This description of $\R$ is convenient but not a foundation: it does not
tell us what an infinite sum is. In \S\ref{sec:axioms} we replace it by a
list of axioms.

\subsubsection{Statements and quantifiers}
A \term{statement} is a sentence that is either true or false. We use the
\term{quantifiers} $\forall$ (``for all'') and $\exists$ (``there
exists''). Each quantifier is written in parentheses, and the statement it
governs in square brackets. For example,
\begin{align*}
    (\forall x\in\R)\,[x^2-x\geq 0] &\quad\text{is false (take } x=\tfrac12\text{)}, \\
    (\exists x\in\R)\,[x^2-x\geq 0] &\quad\text{is true (take } x=2\text{)}.
\end{align*}
Negation interchanges the quantifiers:
\begin{equation*}
    \neg\bigl[(\forall x)\,P(x)\bigr]\iff (\exists x)\,[\neg P(x)],
    \qquad
    \neg\bigl[(\exists x)\,P(x)\bigr]\iff (\forall x)\,[\neg P(x)].
\end{equation*}

The \emph{order} of quantifiers matters, as the next two propositions show.

\begin{proposition}\label{prop:forall-exists}
    $(\forall x\in\R)(\exists y\in\R)\,[x^3-y^3=1]$.
\end{proposition}

\begin{proof}
    Given $x$, take $y=\sqrt[3]{x^3-1}$; then $x^3-y^3=1$.
\end{proof}

\begin{proposition}
    The statement $(\exists y\in\R)(\forall x\in\R)\,[x^3-y^3=1]$ is false.
\end{proposition}

\begin{proof}
    We prove its negation $(\forall y\in\R)(\exists x\in\R)\,[x^3-y^3\neq1]$.
    Given $y$, take $x=y$; then $x^3-y^3=0\neq1$.
\end{proof}

\begin{remark}
    In Proposition~\ref{prop:forall-exists} the number $y$ may depend on
    $x$, and it does: $y=\sqrt[3]{x^3-1}$. The second statement demands a
    single $y$ that works for every $x$ at once, which is much stronger.
    Interchanging $\forall$ and $\exists$ can turn a true statement into a
    false one.
\end{remark}
